\documentclass{article}
\usepackage{tabularx}
\usepackage{booktabs}
\usepackage{geometry}
\geometry{top=1.5cm,left=2cm,bottom=2cm,right=2cm}
\usepackage{enumitem}
\usepackage{ragged2e}

\title{\Large Benchmarking SQL, NoSQL, and Vector Databases for Retrieval-Augmented Generation: A Workload-Specific Performance Analysis}
\author{A.M.A Sakhiur Rahman \\ ID: 2231507642 \and Sadman Jawad Suborno \\ ID: 2233397642}
\date{June 13, 2026}

\begin{document}

\maketitle

\section{Research Question}

\subsection{Primary RQ}
How do relational, document-oriented NoSQL, and purpose-built vector database paradigms differ in retrieval effectiveness and latency across distinct RAG workload types, specifically short-context factual queries, long-context multi-hop queries, and semantic similarity queries, when evaluated against a single domain-specific corpus?

\subsection{Sub-questions}
\begin{itemize}[leftmargin=*]
\item SQ1: Which paradigm achieves the highest precision@10 and NDCG10 per workload type?
\item SQ2: At what query complexity does each paradigm's retrieval latency diverge meaningfully?
\item SQ3: Is there a paradigm that dominates across all three workload types, or do trade-offs consistently emerge?
\end{itemize}

\section{Methodology blueprint}

\subsection{Research paradigm}
Empirical benchmark study (quantitative, positivist).

\begin{table}[h]
\centering
\caption{Systems under study}
\begin{tabularx}{\linewidth}{@{}lXX@{}}
\toprule
Paradigm & System & Rationale \\
\midrule
Relational & PostgreSQL + pgvector extension (via Supabase) & Widely deployed; native vector support added without changing the relational model; represents the "SQL + vector bolt-on" pattern \\
Document NoSQL & MongoDB Atlas (with Vector Search) & Dominant document store with integrated ANN search; represents schema-flexible + vector hybrid \\
Purpose-built vector DB & Qdrant & Designed ground-up for vector workloads; HNSW indexing; represents the "vector-first" paradigm \\
\bottomrule
\end{tabularx}
\end{table}

\subsection{Corpus selection}
ArXiv CS abstracts + full papers (technical, high semantic density, clear ground truth)

\newpage

\subsection{Workload taxonomy}
\begin{table}[h]
\centering
\caption{Workload types in retrieval systems}
\label{tab:workloads}
\begin{tabularx}{\textwidth}{@{}lXXX@{}}
\toprule
Workload type & Definition & Example query & Primary retrieval signal \\
\midrule
W1: Short factual & Single-hop, lookup-style; answer in one chunk & ``What activation function does Transformer use?'' & Keyword + semantic match \\
W2: Multi-hop reasoning & Answer requires combining 2+ chunks across documents & ``How do BERT's training objectives differ from GPT's?'' & Semantic + cross-document linking \\
W3: Semantic similarity & Vague, conceptual queries; no single correct answer & ``Papers about efficient attention mechanisms'' & Pure embedding similarity \\
\bottomrule
\end{tabularx}
\end{table}

\subsection{Metrics}
\begin{enumerate}[leftmargin=*]
\item Primary (retrieval effectiveness):
\begin{itemize}[leftmargin=*]
\item Precision@10, Recall@10
\item NDCG@10
\item Mean Reciprocal Rank (MRR)
\end{itemize}
\item Secondary (operational):
\begin{itemize}[leftmargin=*]
\item Query latency (p50, p95, p99 in ms)
\item Indexing time for the full corpus
\item Memory footprint at rest
\end{itemize}
\end{enumerate}

\subsection{Embedding model}
Fix a single embedding model across all three DBs (e.g., text-embedding-3-smal, text-embedding-3-small from OpenAI, or all-MiniLM-L6-v2 locally) by benchmarking the best one.

\subsection{Experimental controls}
\begin{itemize}[leftmargin=*]
\item Same hardware for all three systems
\item Same chunk size and overlap strategy
\item Same embedding model and dimensionality
\item Each query run 5× and averaged to reduce variance
\item Report standard deviation alongside means
\end{itemize}

\subsection{Analysis approach}
Use paired statistical tests (Wilcoxon signed-rank test) to compare NDCG@10 across paradigms within each workload type

\end{document}